\chapter{Public Network Sniffing}
\label{ch:1.3}
\chaptertopic{Topic 1.3}

\chapterquote{If you know the enemy and know yourself, you need not fear the result of a hundred battles.}{Sun Tzu, \emph{The Art of War}}

\section{Motivation}

Chapters~\ref{ch:1.1} and \ref{ch:1.2} looked at attacks against a human and against a login
form, respectively. This chapter moves one layer down, to the radio medium that carries every
byte you send over Wi-Fi. Public networks -- cafes, airports, hotels -- remove a defense most
people never notice they are relying on: the trust boundary of a private network. At home, you
implicitly trust every device on your Wi-Fi because you (or someone you trust) put it there. On
a public network, every other connected device is a stranger, and the network itself may not be
what it claims to be. This chapter builds a framework for sizing up who might be listening, how
they attack the medium itself, and what a defensible personal protocol for public Wi-Fi looks
like.

\begin{apalignment}{Unit 1: Introduction to Security}{Topic 1.3 -- Best Practices for Public Networks}
\begin{apobjectives}
\apobjective{1.3.A}{Analyze Risk}{Work out what kind of opponent is behind a given attack -- how much capability they have, whether they are reusing tooling somebody else built or writing their own against undisclosed flaws, and what they are trying to get out of it.}
\apobjective{1.3.B}{Analyze Risk}{Recognize the attacks that are specific to radio rather than to cable: a rogue access point impersonating a trusted network, deliberate radio interference that denies service, and roaming surveys that map where a network's signal reaches.}
\apobjective{1.3.C}{Mitigate Risk}{Say what an ordinary user can actually do to keep sensitive data safe over the internet and over Wi-Fi -- checking a network's name character by character before joining, judging whether the data at hand can tolerate an unencrypted link and the metadata such a link exposes, and understanding what a VPN does and does not move out of an adversary's reach.}
\end{apobjectives}
\end{apalignment}

\section{Sizing Up the Adversary}

Not every attacker in a coffee shop is equally dangerous, and treating them all identically
leads either to complacency or to paralysis. Cybersecurity practitioners sort attackers into
tiers based on observable capability, because the tier predicts both how the attack looks and
who is likely to be targeted.

\begin{definition}{Low-skilled adversary}
An attacker (colloquially a ``script kiddie'') who uses existing, publicly available tools
without understanding their underlying mechanics -- downloading a packet-capture app or a
prebuilt evil-twin script rather than writing one.
\end{definition}

\begin{definition}{High-skilled adversary}
An attacker who writes or modifies their own tools, understands the full protocol stack, and can
adapt when defenses push back -- up to and including developing or purchasing
\textbf{zero-day vulnerabilities}: flaws the vendor has not yet discovered, and so has had zero
days to patch.
\end{definition}

The barrier to entry for the low-skilled tier is close to zero -- a thirty-dollar USB Wi-Fi
adapter and a free Linux live-USB are enough to get a novice running a real evil-twin attack.
That low floor is exactly why low-skilled attackers are, in raw numbers, the far more likely
threat for an ordinary student. Millions of people worldwide can run a downloaded script; only a
small population can develop or afford a zero-day. High-skilled attacks are also expensive to
mount, so they are aimed precisely -- executives, journalists, engineers with access to sensitive
repositories. Low-skilled attacks sweep broadly instead, targeting anyone reachable.

\begin{example}[Motivation predicts method]
Attackers do not act at random; their goal shapes their technique. A financially motivated
attacker wants scale -- harvest as many credentials or session cookies as possible per hour, for
resale on underground markets. An espionage-motivated attacker wants one specific thing --
trade secrets, a negotiating position. They will patiently set up a rogue access point near a
single conference room rather than a whole terminal. A disruption-motivated attacker (a
disgruntled former employee, an activist) may jam or flood a network rather than steal anything
at all. Reading an attack's \emph{shape} -- broad and shallow, or narrow and deep -- is often
the fastest way to infer motivation. Motivation, in turn, predicts what the attacker will do
next.
\end{example}

\begin{misconception}
\textbf{``I have nothing worth stealing, so public Wi-Fi is safe for me.''} Low-skilled
attackers do not prescreen victims -- they capture credentials in bulk and resell them without
regard to whose they are. ``Nothing worth stealing'' is a claim about targeted attacks; it says
nothing about opportunistic ones.
\end{misconception}

\section{The Evil Twin Attack}

The most direct exploitation of a public network's missing trust boundary is to impersonate the
network itself.

\begin{definition}{Evil twin}
A rogue access point configured to broadcast the same SSID (network name) as a legitimate
network, often at higher transmit power, so that victim devices connect to the attacker's
hardware instead of the real access point.
\end{definition}

A device does not verify, by default, that the hardware answering to a given network name is the
hardware it connected to last time. An attacker who sets up ``Airport-Free-WiFi'' at higher
signal strength than the real access point wins that competition for most nearby devices.
Devices that auto-connect to remembered network names are especially vulnerable. Once a victim
is on the attacker's hardware, every packet passes through the attacker before it reaches the
internet -- a full \textbf{man-in-the-middle (MitM)} position. Unencrypted HTTP is readable
immediately. HTTPS can sometimes be downgraded with \textbf{SSL stripping}, which tricks a
browser into falling back to plain text. The attacker can also serve a fake
\textbf{captive portal} login page that harvests credentials before quietly forwarding the
victim onward.

\begin{figure}[h]
\centering
\begin{tikzpicture}[
  node/.style={draw=kickstartermaroon, thick, rounded corners, minimum width=2.6cm, minimum height=1cm, align=center, font=\small},
  bad/.style={node, fill=kickstartermaroon!20},
  good/.style={node, fill=kickstartercream}
]
\node[node] (victim) at (0,0) {Victim device};
\node[bad] (evil) at (3.4,1.3) {Rogue AP\\``CafeGuest''\\(higher power)};
\node[good] (real) at (3.4,-1.3) {Real AP\\``CafeGuest''};
\node[node] (internet) at (7,0) {Internet};
\draw[-{Latex[length=2mm]}, kickstartermaroon, thick] (victim) -- (evil) node[midway, above, sloped, font=\scriptsize] {wins auto-connect};
\draw[-{Latex[length=2mm]}, kickstartermaroon, thick] (evil) -- (internet) node[midway, above, sloped, font=\scriptsize] {reads/modifies traffic};
\draw[-{Latex[length=2mm]}, gray, dashed] (victim) -- (real) node[midway, below, sloped, font=\scriptsize, gray] {never reached};
\draw[gray, dashed] (real) -- (internet);
\end{tikzpicture}
\caption{An evil twin broadcasts the same SSID as the legitimate access point, typically at higher
signal strength, so the victim's traffic routes through attacker-controlled hardware -- a full
man-in-the-middle position -- while the real access point never sees the connection at all.}
\label{fig:1.3-evil-twin}
\end{figure}

The tell is a mismatch that the SSID alone cannot reveal. Every physical access point has a
\textbf{BSSID} -- the hardware MAC address of the radio broadcasting it. Two legitimate access
points can share one SSID (normal in an enterprise deployment with many APs), but every BSSID an
organization actually installed is known and approved in advance. An evil twin's BSSID is, by
definition, not on that list. A wireless intrusion detection system (WIDS) automates exactly
this check: maintain the approved BSSID set for each SSID, and alert the moment an unrecognized
BSSID broadcasts an approved name. The same signal generalizes into a short checklist:

\begin{itemize}
\item An approved SSID broadcast from a BSSID outside the known-good set.
\item A \textbf{security downgrade}: the real network runs WPA3, but a second copy of the same
SSID appears at WPA2 or with no security at all. The attacker cannot replicate enterprise
credentials, so a weaker copy is the only option available to them.
\item The same SSID broadcasting under two different security configurations at once.
\item An unexpectedly strong signal from a direction no legitimate AP occupies.
\end{itemize}

\begin{lstlisting}[language=Python, caption={Flagging suspicious access points from a wireless scan}]
def find_evil_twins(access_points, approved_bssids):
    """access_points: list of dicts like
       {"ssid": str, "bssid": str, "security": str}
       approved_bssids: dict mapping ssid -> set of known-good BSSIDs.
       Returns the sorted set of BSSIDs that are suspicious because their
       SSID is claimed by more than one BSSID, or because the BSSID is not
       in the approved set for that SSID."""
    by_ssid = {}
    for ap in access_points:
        by_ssid.setdefault(ap["ssid"], []).append(ap)

    suspicious = set()
    for ssid, aps in by_ssid.items():
        bssids_seen = {ap["bssid"] for ap in aps}
        approved = approved_bssids.get(ssid, set())
        if len(bssids_seen) > 1:
            suspicious.update(bssids_seen - approved)
        for ap in aps:
            if ap["bssid"] not in approved:
                suspicious.add(ap["bssid"])
    return sorted(suspicious)
\end{lstlisting}

\begin{misconception}
\textbf{``If I see HTTPS in my browser, the connection is safe even on public Wi-Fi.''}
SSL stripping can remove HTTPS before the page ever reaches your browser. Its absence where you
expected it -- a site that has always been secure suddenly is not -- is itself a warning sign,
not a neutral fact.
\end{misconception}

\section{Denial of Service by Radio Jamming}

An attacker does not need to intercept traffic to hurt you; making the network unusable at all
is itself an attack.

\begin{definition}{Radio-frequency (RF) jamming}
Deliberate transmission of noise on a wireless channel to prevent legitimate devices from
communicating, exploiting the channel-sensing behavior every Wi-Fi radio depends on.
\end{definition}

Wi-Fi devices share spectrum politely using \textbf{CSMA/CA} (Carrier Sense Multiple Access with
Collision Avoidance). Before transmitting, a radio listens, and if it hears noise -- legitimate
traffic or otherwise -- it backs off and waits. A jammer weaponizes that courtesy: by emitting
noise continuously above the ambient level, it forces every legitimate device nearby to defer
indefinitely, at almost no cost to the attacker. The effect is a \textbf{denial-of-service (DoS)}
-- the network is unusable not because anyone was authenticated out of it, but because the
medium itself has been rendered noisy.

The practical challenge is distinguishing an attack from ordinary interference. Microwave ovens,
Bluetooth devices, and neighboring Wi-Fi networks all raise the 2.4~GHz noise floor briefly and
constantly; that is normal background texture, not an attack. What separates jamming is
\emph{duration}: a sustained, elevated noise plateau across many consecutive measurements.
Attackers often time it to coincide with something they want disrupted -- an exam, a hearing, a
physical intrusion where wireless alarms need to go quiet.

\begin{lstlisting}[language=Python, caption={Detecting a sustained jamming event in a noise-floor sample series}]
def is_jamming(samples, threshold, min_consecutive):
    """samples: list of int noise-floor readings.
       Returns True iff at least `min_consecutive` consecutive samples
       exceed `threshold`."""
    run = 0
    for value in samples:
        run = run + 1 if value > threshold else 0
        if run >= min_consecutive:
            return True
    return False
\end{lstlisting}

Jamming devices are illegal for civilian use in the United States under 47 U.S.C. \S 333, but
illegality does not stop use in targeted disruption of conference venues or exam halls, or as
cover for a physical break-in against wireless alarms and cameras. Defenders cannot always
prevent jamming outright, so detection is usually the actionable response -- the same ``sustained
anomaly over a window'' logic used throughout intrusion detection. Two complementary steps help:
switching to a less-congested band (5~GHz or 6~GHz have more channels and shorter range than
2.4~GHz), and alerting venue security.

\begin{misconception}
\textbf{``Switching Wi-Fi channels defeats a jammer.''} It helps against a jammer targeting one
specific channel, but a wideband jammer can cover an entire 2.4~GHz band simultaneously,
making channel-hopping ineffective on its own.
\end{misconception}

\section{War Driving and the Personal Protocol}

\begin{definition}{War driving}
Moving through a geographic area with a wireless scanner to discover, log, and map Wi-Fi
networks -- their SSIDs, BSSIDs, security types, and locations.
\end{definition}

The term dates to the early 2000s, when attackers physically drove through business districts
with directional antennas hunting for unsecured corporate networks. Today a smartphone scanning
app does the same job. War-driving data has legitimate uses -- coverage planning, academic
research -- and malicious ones: identifying open networks, networks still running long-broken
security (WEP was cracked in 2005), and networks whose SSID leaks an organization's identity. In
the worst case, all of this is reachable from a public parking lot with no need to set foot
inside a building.

As a user rather than a scanner, you face the mirror image of this same risk when you open your
device's Wi-Fi list in public. Three properties of a network compound into real exposure:

\begin{itemize}
\item \textbf{Open network.} No password means no encryption at the Wi-Fi layer -- anyone
sharing the channel can passively capture your traffic.
\item \textbf{Unverified SSID.} A plausible-sounding name (``AirportFreeWifi'') is not proof of
anything; it may be an evil twin from earlier in this chapter.
\item \textbf{No VPN.} Even on a legitimate, encrypted network, the access-point operator can
observe your DNS queries and any unencrypted traffic; on an open network, any co-joined device
can too.
\end{itemize}

These three risks resolve into an ordered protocol, and the order matters. \textbf{(1) Verify the
SSID} with venue staff before connecting -- connecting to the wrong network at all makes every
later protection moot. \textbf{(2) Avoid open networks} for sensitive activity, since even a few
unencrypted seconds before a VPN tunnel initializes can leak a credential. \textbf{(3) Always use
a VPN}, which encrypts traffic to a trusted exit point regardless of what the local network does
-- the defense-in-depth layer that still helps when the first two steps are imperfect.

\begin{lstlisting}[language=Python, caption={Flagging the risk properties of a network before connecting}]
def assess_network(network):
    """network: dict like {"ssid": str, "security": str,
       "verified": bool, "has_vpn": bool}
       Returns a sorted list of risk-flag strings that apply."""
    flags = []
    if network["security"] in ("none", "open"):
        flags.append("open-network")
    if not network["verified"]:
        flags.append("unverified-ssid")
    if not network["has_vpn"]:
        flags.append("no-vpn")
    return sorted(flags)
\end{lstlisting}

\begin{example}[Reconnaissance feeding social engineering]
War-driving data does not only enable direct wireless attacks. An attacker who has logged your
organization's SSID from a public war-driving database can fold that fact into a phishing email
-- ``Your corporate Wi-Fi password expires tomorrow, click here to renew.'' The result carries far
more plausibility than a generic lure. Physical-layer reconnaissance and the social-engineering
techniques from Chapter~\ref{ch:1.1} are not separate problems; sophisticated attacks routinely
chain both together.
\end{example}

\begin{misconception}
\textbf{``A password-protected network means I'm safe.''} WPA2 or WPA3 encrypts the air
interface between your device and the access point. But the AP operator -- or an evil twin
posing as one -- can still read whatever crosses that boundary in the clear. A VPN protects the
layer that Wi-Fi security does not.
\end{misconception}

\section{Case Studies in Wireless Compromise}

\begin{example}[DarkHotel: A Targeted, Not Opportunistic, Evil Twin]
Beginning around 2010 and characterized publicly by Kaspersky in 2014, the DarkHotel campaign
compromised the networks of luxury hotels and waited for a specific, named guest -- an executive
or senior technical staffer in defense, pharmaceuticals, or private equity -- to check
in.\srccite{1.3}{1} Rather than harvesting whatever traffic passed, the operators served that one
guest -- and no one else on the same network -- a fake software-update prompt, mimicking Adobe
Flash or a browser toolbar. It installed a surveillance implant alongside a working copy of the
real software. The prompt was convincing because the certificates behind it were cryptographically
forged. The operators had factored the weak keys of real 512-bit RSA certificates and used the
recovered private keys to sign their malicious installer, so it validated as genuine. This case
illustrates two points beyond the evil-twin mechanics already covered. First, an attack can be
surgically targeted at one occupant of a shared network rather than broadcast to everyone on it,
evading detection strategies built around bulk anomalies. Second, the abused trust here was the
software-update channel, not a login page. That is exactly why operating systems now verify update
signatures against a pinned publisher key independent of the network that delivered them.
\end{example}

\begin{example}[KRACK, 2017: A Fixed Protocol, an Unfixed Fleet]
In October 2017 researcher Mathy Vanhoef disclosed KRACK (Key Reinstallation Attack), a flaw in how
the WPA2 handshake itself handled retransmitted messages.\srccite{1.3}{2} Nearly every WPA2 device
shipped since 2004 was technically vulnerable as a result, independent of password strength. The
fix was a single rule (never reinstall a key already held) and was simple enough that major
operating systems patched within days of disclosure. What makes KRACK a case study rather than a
footnote is what happened next. Laptops, phones, and actively maintained software received the
patch almost immediately. Consumer routers past their support window, unsupported Android
handsets, and embedded devices with firmware burned in at manufacture never did -- and, by design,
never could. The population of exploitable devices shrank only as those devices were discarded,
not as they were fixed. Vulnerability severity and real-world exposure are governed by different
variables. A protocol flaw is only as fixable as the update pathway of the hardware it runs on.
That is why ``how would a fix reach every deployed unit, and who maintains that pathway'' belongs
in a system's initial design, not only in its incident response plan.
\end{example}

\begin{misconception}[``A prosecuted evil twin case means these attacks are rare in practice.'']
An Australian man arrested in April 2024 ran a portable evil-twin access point across three
airports and multiple domestic flights; he was ultimately sentenced to over seven years. The case
shows the opposite of rarity. The seized equipment was a commercially sold, off-the-shelf portable
wireless auditing device requiring no exploit development at all, and the fraudulent login pages
found on it requested nothing more sophisticated than an email or social-media login ``to
connect.'' Evil-twin attacks are common precisely because commodity hardware and a fake login page
are the entire attack. No vulnerability in Wi-Fi security itself is required, since the operator
never needs to decrypt HTTPS traffic when a captive portal can simply ask the victim to type a
credential directly into an attacker-controlled page. That a legitimate network never needs a
password or a social-media login just to grant an IP address is, on its own, the single most
reliable tell in this chapter.
\end{misconception}

\section*{Review Questions}
\begin{reviewquestions}
\item A wireless scan shows the SSID \texttt{"CafeGuest"} broadcast from three different BSSIDs,
only two of which appear on the venue's approved list. Explain what conclusion this supports and
why BSSID, not SSID, is the reliable signal.
\item Explain why low-skilled adversaries are statistically the more likely threat to an
ordinary user on public Wi-Fi, even though high-skilled adversaries are far more dangerous in
absolute terms.
\item A noise-floor monitor logs three seconds of elevated readings, then returns to baseline, then
logs three minutes of elevated readings. Explain which event is more likely a jamming attack and
what property of the data distinguishes it from ordinary interference.
\item Explain why the personal protocol for public Wi-Fi (verify SSID, avoid open networks, use a
VPN) is presented in that specific order rather than any other, referencing what each earlier
step protects against that a later step cannot.
\item A network is open, has an SSID you have not verified with staff, and you have no VPN
running. Using the risk-flagging logic from this chapter, list the flags this network would
raise and explain the compounding risk of connecting anyway.
\item In the DarkHotel campaign, the attackers served a malicious update prompt to one specific
guest and no one else on the same hotel network. Explain why this targeting choice made the
attack harder for bulk-anomaly-based detection to catch, compared to an evil twin that harvests
every connected device's traffic indiscriminately.
\item KRACK was patched within days on laptops and phones but remained exploitable for years on
many consumer routers and embedded devices. Explain why the same disclosed vulnerability produced
such different real-world exposure across device types, and what this implies about evaluating a
device's security before purchase.
\end{reviewquestions}

\begin{answerkey}
\item Three BSSIDs answering to one SSID, with only two on the approved list, supports the
conclusion that a rogue access point (an evil twin) is impersonating the legitimate network. SSID
is reliable evidence of nothing because any device can broadcast any name it chooses. BSSID
identifies the specific radio hardware instead, and an organization's approved BSSID set is fixed
and known in advance. An unlisted BSSID broadcasting a trusted name is, by definition, hardware
the organization did not install.
\answerlabel{1.3.B}
\item Low-skilled adversaries require only inexpensive, widely available tools and no specialized
knowledge, so the population capable of mounting such an attack is enormous, and they sweep broadly
across whatever targets are reachable. High-skilled adversaries are far more dangerous per
incident, but developing or acquiring the capability (custom tooling, zero-days) is expensive
enough that they are deployed sparingly and aimed at specific, high-value targets. For an ordinary
user, the sheer numerical prevalence of low-skilled attackers makes them the statistically more
likely threat even though any single high-skilled attack could do more damage.
\answerlabel{1.3.A}
\item The three-minute sustained elevated reading is more likely a jamming attack; the three-second
spike is more consistent with ordinary interference such as a microwave oven or a neighboring
device briefly raising the noise floor. What distinguishes an attack is duration: a jammer must
continuously emit noise to keep legitimate devices deferring, so a sustained plateau across many
consecutive measurements is the signature that ordinary, transient interference does not produce.
\answerlabel{1.3.B}
\item The order matters because each step protects against a failure the later steps cannot
recover from. Verifying the SSID comes first because connecting to the wrong network at all makes
every subsequent protection irrelevant -- an evil twin can still present a working VPN-compatible
connection while intercepting the negotiation. Avoiding open networks comes second because even a
brief window of unencrypted traffic before a VPN tunnel initializes can leak a credential that no
later encryption can retroactively protect. Using a VPN comes last as the defense-in-depth layer
that still helps even when the first two steps are imperfectly executed, encrypting traffic to a
trusted exit point regardless of what the local network or access point operator can see.
\answerlabel{1.3.C}
\item This network raises all three flags. Open-network, because no password means no Wi-Fi-layer
encryption. Unverified-ssid, because a plausible name has not been confirmed with venue staff and
may belong to an evil twin. No-vpn, because nothing encrypts traffic to a trusted exit point. The
risks compound rather than add: an unverified network might itself be hostile hardware, and an open
configuration means any nearby device can passively capture whatever crosses it. The absence of a
VPN removes the one layer that would still protect the connection even if the first two assumptions
turned out to be wrong.
\answerlabel{1.3.C}
\item Bulk-anomaly detection looks for patterns like many devices suddenly receiving the same
unexpected prompt, or traffic volumes spiking network-wide -- signals that a broadcast evil twin
naturally produces. By serving the malicious prompt to exactly one guest, DarkHotel generated no
network-wide anomaly at all: everyone else's traffic looked completely normal, so any detection
strategy built around aggregate deviation had nothing unusual to flag. Only monitoring focused on
that individual guest's device, or verifying update signatures independent of the network, could
have caught it.
\answerlabel{1.3.A}
\item The vulnerability was identical everywhere, but exposure was governed by whether each
device had a maintained update pathway. Laptops and phones run actively supported operating
systems with vendors who ship patches quickly. Many consumer routers and embedded devices, by
contrast, either had no remaining support window or had firmware burned in at manufacture with no
update mechanism at all. This implies that evaluating a device's security before purchase should
include its vendor's patch-support timeline and update mechanism, not just its security features at
the moment of sale. A device's real-world security depends on how long, and how reliably, that
pathway will remain open.
\answerlabel{1.3.B}
\end{answerkey}

\section*{Sources}
\begin{enumerate}[label=\lbrack\arabic*\rbrack, leftmargin=2.2em]
\item \srcitem{1.3}{1}Kaspersky Lab, Global Research and Analysis Team, \emph{The Darkhotel APT},
version 1.1, November 2014.
\item \srcitem{1.3}{2}Mathy Vanhoef and Frank Piessens, ``Key Reinstallation Attacks: Forcing Nonce
Reuse in WPA2,'' \emph{Proceedings of the 2017 ACM SIGSAC Conference on Computer and Communications
Security (CCS '17)}, October 2017; publicly disclosed October 16, 2017.
\end{enumerate}
